\section{9. Discussion}

\subsection{9.1 Key Takeaways}

\subsubsection{9.1.1 Takeaway 1: Linear Attention Without Random Features}

ClusterFormer achieves linear complexity \textbf{without approximating softmax}. 
It's exact on the cluster level, differentiating it from Performer and 
Linformer.

\subsubsection{9.1.2 Takeaway 2: Global + Local is Optimal}

Neither global clustering alone nor local window alone matches the blend. 
The learnable gate discovers the optimal mix per layer.

\subsubsection{9.1.3 Takeaway 3: Not a Compression Method}

Parameter count increases by LoRA overhead. ClusterFormer targets 
\textbf{compute efficiency}, not size reduction.

\subsection{9.2 When to Use ClusterFormer}

\subsubsection{Recommended}

\begin{itemize}
\item \textbf{Long sequences} ($L \geq 256$)
\item \textbf{Attention-bound workloads}
\item \textbf{Research on clustering attention}
\item \textbf{Multi-architecture pipelines}
\end{itemize}

\subsubsection{Not Recommended}

\begin{itemize}
\item \textbf{Short sequences} ($L < 100$) -- overhead exceeds savings
\item \textbf{Model size reduction} -- use compression methods
\item \textbf{Strict causal generation} -- decoder uses exact attention anyway
\item \textbf{Interpretability-first} -- attention maps are cluster-level
\end{itemize}

\subsection{9.3 Decision Flowchart}

\begin{lstlisting}
Need size reduction? --- Yes --- Use A3/AASVD/Adaptive SVD
  | No
   
Sequence length > 128? --- No --- Use baseline softmax
  | Yes
   
Latency critical? --- Yes --- Consider FlashAttention
  | No
   
Use ClusterFormer
\end{lstlisting}

\subsection{9.4 Cost-Benefit}

For DistilBERT at $L=256$:

\begin{table}[H]
\centering
\small
\begin{tabular}{llll}
\toprule
\textbf{Metric} & \textbf{Baseline} & \textbf{ClusterFormer} & \textbf{Change} \\
\midrule
Attention FLOPs & 6.7M & 2.4M & **2.8$\times$ reduction** \\
Accuracy & 90.71\% & 89.91\% & -0.80\% \\
Params & 66.9M & 67.55M & +589826M \\
Wall time & 4.2 ms & 3.8 ms & 1.10$\times$ \\
\bottomrule
\end{tabular}
\end{table}

\subsection{9.5 Comparison with Published Methods}

\begin{table}[H]
\centering
\small
\begin{tabular}{lllll}
\toprule
\textbf{Method} & \textbf{Type} & \textbf{Complexity} & \textbf{Global} & \textbf{Local} \\
\midrule
Longformer & Sparse & $O(LW)$ & Partial & Yes \\
Performer & Kernel & $O(Lm)$ & Yes & No \\
Reformer & LSH & $O(L\log L)$ & Yes & No \\
**ClusterFormer** & **Cluster** & **$O(LK)$** & **Yes** & **Yes** \\
\bottomrule
\end{tabular}
\end{table}

\subsection{9.6 Deployment Considerations}

\subsubsection{Latency}

Beneficial for $L \geq 128$. Below that, overhead dominates.

\subsubsection{Memory}

Attention matrix is $O(KL)$ instead of $O(L^2)$: \textbf{substantial savings} 
for long sequences.

\subsubsection{Interpretability}

Cluster assignments give coarse interpretability; per-token attention is 
not directly accessible.

\subsection{9.7 Open Questions}

1. \textbf{Optimal $K$ per layer} -- could be learned dynamically?
2. \textbf{Interaction with quantization} -- INT8 + ClusterFormer?
3. \textbf{Causal decoding} -- can clustering work for generation?
4. \textbf{Very long context} ($L > 4096$) -- does the crossover hold?

\subsection{9.8 Broader Impact}

\subsubsection{Positive}
\begin{itemize}
\item Compute reduction for long-context inference
\item Democratizes long-context models
\item Complements other efficiency methods
\end{itemize}

\subsubsection{Negative}
\begin{itemize}
\item Modest parameter increase
\item Clustering opacity
\item Dual-use
\end{itemize}

\subsubsection{Environmental}
Total compute: ~15 GPU-hours, ~{{CO2}} kg CO$_2$eq.